\section*{Bangladesh vs World: Dependency Ratios (CDR, ADR, TDR)}

\noindent\textbf{What the ratios mean.}
\begin{itemize}
  \item \textbf{Child Dependency Ratio (CDR)}: dependents aged 0--14 per 100 people aged 15--64.
  \item \textbf{Aged (Old-age) Dependency Ratio (ADR)}: dependents aged 65+ per 100 people aged 15--64.
  \item \textbf{Total Dependency Ratio (TDR)}: CDR + ADR; the overall burden on the working-age population.
\end{itemize}
\begin{figure}[H]
    \centering
    \includegraphics[width=0.5\linewidth]{Picture1.png}
     
    \label{fig:placeholder}
\end{figure}
\begin{figure}[H]
    \centering
    \includegraphics[width=0.5\linewidth]{Picture2.png}
     
    \label{fig:placeholder}
\end{figure}
\begin{figure}[H]
    \centering
    \includegraphics[width=0.5\linewidth]{Picture3.png}
    
    \label{fig:placeholder}
\end{figure}

\noindent\textbf{How to read the chart.}
\begin{itemize}
  \item Each line shows the ratio through time. A \emph{declining CDR} usually signals falling fertility and progress through the demographic transition.
  \item A \emph{rising ADR} indicates population ageing; this can lift the TDR even when the CDR is falling.
  \item The \emph{gap between Bangladesh and the World} tells you whether Bangladesh’s pressure on workers is above or below the global average, and whether it is converging.
\end{itemize}

 

\paragraph{Key comparison (typical pattern seen in the graph).}
\begin{itemize}
  \item \textbf{CDR:} Bangladesh generally starts above the World average (younger age structure) and \emph{declines faster}, shrinking the gap over time. This reflects sustained fertility decline.
  \item \textbf{ADR:} Bangladesh’s ADR is typically \emph{below} the World average for many years but \emph{rises steadily} as longevity improves and cohorts age. The gap tends to narrow in the most recent years.
  \item \textbf{TDR:} With CDR falling faster than ADR rises, Bangladesh’s TDR shows a clear \emph{downward trend}, indicating a lighter overall burden on workers and an extended \emph{demographic dividend} window. In recent years, TDR approaches (or converges toward) the World average.
\end{itemize}

\paragraph{Implication.}
Bangladesh’s rapid CDR decline has expanded its working-age share, supporting growth and savings. However, the gradual rise in ADR means policies for healthy ageing, pensions, and productivity gains will matter more going forward, so that TDR remains manageable as the population structure matures.
\section{Reference}



\end{document}

